\section{Sovereign Business Connector Onboarding}
\label{sec:sovereign-business-connectors}

This section records completion of roadmap task \texttt{SBB-0207}.  The
implementation does not create a new Files, mail, calendar, sales or accounting
system.  It presents the connector authorities already present in AION and the
Boardroom through one owner-facing onboarding surface and projects their safe
status into the existing sovereign business architecture.

\subsection{Product behaviour}

The Boardroom now contains a single \emph{Connect your business} area with five
plain-language cards:

\begin{itemize}
  \item \textbf{Files} opens the existing governed File Cabinet;
  \item \textbf{Gmail} opens the established Google authorization route;
  \item \textbf{Google Calendar} reuses the persona-bound Google authorization
        while checking separately for calendar scope;
  \item \textbf{CRM/Sales} reuses the existing HubSpot connection authority and
        the Sales Revenue Spine; and
  \item \textbf{Xero} reuses the existing customer-business Xero authorization
        and Business Financial Model boundary.
\end{itemize}

Each card answers only three immediate questions: whether the source is connected,
whether its local connector is healthy, and how much evidence is already represented
in the appropriate business container.  The Boardroom remains usable when no external
account is connected.  Failure of one provider is isolated to its card and does not
hide local files or healthy services.

\subsection{Owner initiation and absence of side effects}

Loading or refreshing the onboarding screen is a read-only operation.  It cannot
create an account, start OAuth, import mail, alter a calendar, synchronize accounting
records or modify the Business Map.  An authorization flow begins only after the
owner presses the relevant connection button.  The provider authorization page is
opened through the existing desktop external-browser boundary, and the owner returns
to refresh the connection status after granting permission.

The screen is consequently a connection and evidence-status surface, not an implicit
consent mechanism.  A connected state proves that a credential reference and required
scope are available to the mother brain; it does not prove that every provider item
has been imported, reviewed or accepted as company truth.

\subsection{Architecture and data flow}

The implementation introduces a provider-neutral
\texttt{BusinessConnectorOnboardingService}.  It composes five existing authorities
rather than duplicating them.  Its output is exposed through one local endpoint and
rendered in the existing Boardroom Business Map governance panel.

\begin{verbatim}
Owner opens Boardroom
        |
        v
Read-only connector onboarding service
        |
        +-- File Cabinet ----------> business knowledge / files
        +-- Gmail ----------------> Sales Revenue Spine
        +-- Calendar -------------> calendar planning + appointment evidence
        +-- CRM / HubSpot --------> Sales Revenue Spine + Sales Intelligence
        +-- Xero -----------------> Business Financial Model
        |
        v
Secret-free health + aggregate projection counts
        |
        v
Existing review boundary -> canonical Business Map (only if approved)
\end{verbatim}

No connector receives direct write authority over the canonical Business Map.
Provider records first remain source evidence.  Existing provenance, confidence,
review-state, correction, retention, revocation and legal-hold controls determine
whether a derived claim may later become canonical business knowledge.

\subsection{Projection contract}

The unified response uses schema
\texttt{aion.business\_connector\_onboarding.v1}.  Every connector card contains a
stable identifier, a non-technical label, connected and health states, its target
business container, an aggregate projection summary, and the existing owner-initiated
connection route.  It also declares that onboarding is read-only and that raw provider
data is not Business Map truth.

The current projections are deliberately narrow:

\begin{itemize}
  \item Files reports document and protected-record counts;
  \item Gmail reports the count of email-origin enquiries already represented in the
        Sales Revenue Spine;
  \item Calendar reports the count of appointment signals already represented;
  \item CRM/Sales reports CRM-origin opportunity counts and aggregate pipeline stages;
  \item Xero reports bounded synchronization status and its existing safe summary.
\end{itemize}

These are navigation and readiness summaries.  They are not financial advice,
delivery proof, calendar truth or a claim that every external record has been mapped.

\subsection{Privacy and credential boundary}

Provider secrets remain inside the customer-controlled mother brain and existing
vault authorities.  The onboarding response and its user interface exclude passwords,
access tokens, private keys, raw messages, contact details, calendar event titles and
bank details.  Email and CRM projections count records without returning customer names
or addresses.  Google is considered connected only when the existing connector reports
a connected authorization and a real token reference; Calendar additionally requires
an authorized calendar scope.

This separation permits the same interface to work in local, customer-cloud and
enterprise deployments without making the desktop page a credential store.  A future
provider can be added behind the same card contract while its authentication, scopes,
residency and data handling remain governed by its own adapter.

\subsection{Failure and trust semantics}

Connector probing is fail-closed and fault-isolated.  If a provider cannot be reached,
its card reports unavailable rather than connected.  Other cards still render.  Missing
Google token possession or missing Calendar scope produces a not-connected state even
if a superficial provider status says connected.  The interface never converts a
transport response into evidence acceptance, nor a successful OAuth page into a claim
that synchronization completed.

The visual language intentionally distinguishes \emph{connected}, \emph{not connected}
and evidence already \emph{represented}.  Further states---such as imported, reviewed,
approved, rejected, stale and revoked---remain responsibilities of the underlying
container and Business Map governance layers.

\subsection{Verification and release record}

Automated tests cover all five cards, Google credential and Calendar-scope checks,
aggregate Files, Gmail, Calendar and CRM/Sales projections, Xero handoff metadata,
absence of private source fields, provider-failure isolation, API error handling and
the existing Boardroom controls.  Static interface tests confirm that account routes
are invoked only from explicit button handling and that Files reuses the existing File
Cabinet surface.

The capability is packaged in AION Sovereign Brain Bootstrap~0.5.0 and Pilot
Fabric~0.57.0.  The developer-preview installers remain unsigned; production signing,
clean-machine qualification and real-account field tests remain separate release gates.

\subsection{Phase closure and next work}

This completes the locally implementable tasks in Phase~2 of the Sovereign Business
Brain roadmap: the original Boardroom Business Map remains canonical; onboarding is
resumable; facts and relationships are governed; review queues and owner editing are
available; approved source claims can be linked deterministically; department access is
bounded; source retention and revocation are enforced; and the established business
connectors now have one safe onboarding view.

The next implementation phase is the unified intelligence contract.  It must define one
provider-neutral request, structured response and streaming envelope across deterministic
capabilities, retrieval, local models, customer-cloud models and optional external APIs,
while preserving explicit disclosure, residency, budget, fallback and execution limits.
